\documentclass[11pt,a4paper]{article}
\usepackage[margin=1in]{geometry}
\usepackage{times}
\usepackage{latexsym}
\usepackage{amsmath}
\usepackage{amssymb}
\usepackage{booktabs}
\usepackage{graphicx}
\usepackage{multirow}
\usepackage{xcolor}
\usepackage{url}
\usepackage{natbib}
\usepackage[colorlinks=true,citecolor=blue,urlcolor=blue]{hyperref}

\title{The Depth~2 Ceiling: Generational Signals in Recursive Synthetic Text\\Collapse Beyond First-Generation Detection}

\author{Anonymous}

\begin{document}
\maketitle

\begin{abstract}
Web-scale corpora are increasingly contaminated by recursive synthetic text---machine-generated content that is itself generated from prior machine outputs.
While binary detection (human vs.\ synthetic) is well studied, estimating the \emph{generational depth} of such text remains unexplored.
We formalize recursive generation as an iterated lossy channel and derive from the Data Processing Inequality (DPI) that a \emph{saturation boundary} $K^*$ must exist beyond which consecutive generations become statistically indistinguishable.
Through controlled experiments across three model families (Qwen2.5-1.5B, Pythia-1.4B, OLMo-1B), four decoding strategies, and a 7B scale probe, we establish $K^*=2$ as a hard, universal ceiling for models in the 1--7B range: pairwise Jensen--Shannon divergence decays from 0.27 (depth 0 vs.\ 1) to $<$0.005 (depth 2 vs.\ 3), and no combination of model architecture, decoding strategy, or feature representation shifts this boundary.
Consequently, graduated depth-aware filtering offers no statistically significant advantage over binary filtering on downstream benchmarks ($\Delta\leq$0.2pp, $p>0.4$).
These negative results carry direct practical implications: depth-aware data curation has an extremely narrow actionable window (human / gen-1 / gen-2+), and community efforts should focus on first-generation detection rather than fine-grained depth estimation.
\end{abstract}

\section{Introduction}

The proliferation of large language models (LLMs) has led to a feedback loop in which machine-generated text enters web corpora, is scraped, and is used to train subsequent models~\cite{shumailov2024ai}.
This \emph{recursive contamination} raises data quality concerns that binary synthetic-text detectors~\cite{mitchell2023detectgpt,hans2024binoculars} are not designed to address: they classify text as human or machine-generated but do not distinguish first-generation synthetic text (relatively high quality) from fourth-generation text (potentially degenerate).

A natural question arises: \emph{can we estimate the generational depth of synthetic text, and does such estimation improve data curation?}
In this work, we provide a rigorous negative answer to both questions for current-scale (1--7B) language models.

We model recursive text generation as an iterated lossy channel, where each generation pass loses mutual information with the original human distribution.
The Data Processing Inequality (DPI) guarantees the existence of a \emph{saturation boundary} $K^*$---the maximum depth at which consecutive generations remain distinguishable.
Our contribution is the first systematic measurement of $K^*$ across controlled experimental conditions, yielding a stark finding:

\begin{itemize}
    \item $K^*=2$ is invariant across three model families (Qwen2.5, Pythia, OLMo), four decoding strategies (greedy, nucleus $p$=0.9, $T$=0.7, $T$=1.2), base vs.\ instruct variants, and scales from 1B to 7B parameters.
    \item Graduated filtering (exploiting depth beyond the binary distinction) provides $\leq$0.2 percentage points of improvement over binary filtering, statistically indistinguishable from zero ($p>0.4$).
    \item Neither learned representations (frozen LM embeddings) nor frontier LLMs (GPT-4o, Claude~3.5) can penetrate the $K^*=2$ ceiling.
\end{itemize}

These results constitute a \emph{cautionary message} for the community: the window for depth-aware data curation is extremely narrow---effectively a three-class problem (human / gen-1 / gen-2+)---and the marginal value of depth estimation over binary detection is negligible at current model scales.
Data governance efforts are better directed at robust first-generation detection.

\section{Background and Theoretical Framework}

\subsection{Recursive Generation as an Iterated Lossy Channel}

Let $P_0$ denote the human text distribution and $P_k$ the distribution after $k$ rounds of generation through a language model $\mathcal{M}$.
Each generation step applies a stochastic mapping $T_\mathcal{M}: P_{k-1} \mapsto P_k$.
We model $T_\mathcal{M}$ as a lossy channel: while $\mathcal{M}$ approximates its input distribution, it introduces systematic biases (lower perplexity variance, reduced lexical diversity, peaked token distributions) that constitute information loss.

\subsection{Data Processing Inequality and the Saturation Boundary}

By the DPI, for any Markov chain $P_0 \to P_1 \to \cdots \to P_k$:
\begin{equation}
    I(P_0; P_k) \leq I(P_0; P_{k-1}) \leq \cdots \leq I(P_0; P_1)
\end{equation}
where $I(\cdot;\cdot)$ denotes mutual information.
The Jensen--Shannon divergence between consecutive generations $\text{JSD}(P_{k-1}, P_k)$ is bounded above by the mutual information, and must therefore be monotonically non-increasing.
We define the \emph{saturation boundary}:
\begin{equation}
    K^* = \max\{k : \text{AUC}(P_{k-1}, P_k) > \tau\}
\end{equation}
where $\text{AUC}(P_{k-1}, P_k)$ is the area under the ROC curve for pairwise binary classification between generations $k{-}1$ and $k$, and $\tau$ is a discrimination threshold (default $\tau=0.60$).

The DPI guarantees that $K^*$ is finite for any lossy channel.
Our empirical question is: \emph{how small is $K^*$ in practice?}

\subsection{Related Work}

\paragraph{Model Collapse.} \citet{shumailov2024ai} demonstrate that iterative training on model outputs leads to distributional collapse. Our work measures the \emph{speed} of this collapse in the feature space relevant to detection, rather than its long-run consequences for model quality.

\paragraph{Binary Synthetic Text Detection.} Methods such as DetectGPT~\cite{mitchell2023detectgpt}, Binoculars~\cite{hans2024binoculars}, and DivEye treat detection as a binary task. We show that the confidence scores of these detectors carry negligible depth information (Spearman $\rho \leq 0.22$; depth-2+ $\rho < 0.05$).

\paragraph{Data Quality Filtering.} Verification-based approaches~\cite{feng2024towards,yi2024on} use external verifiers for binary quality gating. We investigate whether \emph{ordinal} depth signals can refine such filtering, and find that they cannot.

\section{Experimental Setup}

\subsection{Depth Chain Construction}

For each model $\mathcal{M}$ and seed corpus, we construct a \emph{depth chain}: a sequence of corpora $\{D_0, D_1, \ldots, D_5\}$ where $D_0$ is human text and $D_k$ ($k \geq 1$) is generated by prompting $\mathcal{M}$ to continue or paraphrase documents from $D_{k-1}$.
We use 5,000 documents from C4 as the seed corpus $D_0$.

\paragraph{Models.} Qwen2.5-1.5B-Instruct (primary), Pythia-1.4B, OLMo-1B, and Qwen2.5-7B (scale probe, 1K documents, depth 0--3).

\paragraph{Decoding Strategies.} Nucleus $p$=0.95 $T$=1.0 (default), greedy, nucleus $p$=0.9, $T$=0.7, and $T$=1.2.

\subsection{Feature Extraction}

We extract 15 handcrafted features organized into three groups:
\textbf{Surprisal statistics} (mean, variance, skewness of token-level log-probabilities under an external scorer),
\textbf{Diversity metrics} (type-token ratio, hapax ratio, entropy of unigram distribution), and
\textbf{Tail-thinning indicators} (fraction of tokens in the bottom-5\% probability mass, kurtosis of the token distribution).

As a control, we also evaluate \textbf{learned representations}: frozen Qwen2.5-1.5B last-layer embeddings (4096-dim) with a 2-layer MLP.

\subsection{Depth Estimation Methods}

We compare four approaches: (1) \textbf{CORAL}~\cite{cao2020rank}---ordinal classification with consistent rank predictions; (2) \textbf{Standard multiclass}---softmax cross-entropy, no ordinal structure; (3) \textbf{Ordinal regression}---continuous depth prediction with discretization; (4) \textbf{Binary + random depth}---binary (human vs.\ synthetic) classifier followed by uniform random depth assignment for synthetic texts.

\subsection{Downstream Evaluation}

We evaluate four data filtering strategies via continued pretraining of Qwen2.5-1.5B on filtered corpora, followed by evaluation on MMLU and HellaSwag (3 seeds):
\textbf{No filter}, \textbf{Binary} (remove all predicted synthetic), \textbf{Graduated} (depth-weighted inclusion), and \textbf{Surprisal-based} (continuous surprisal score as quality proxy).

\section{Results}

\subsection{$K^*=2$ Is a Universal Ceiling}

Table~\ref{tab:cross_model} presents pairwise discrimination results across three model families.
All models yield $K^*=2$ at $\tau=0.60$: the 0-vs-1 boundary is nearly perfectly separable (AUC $\geq$ 0.988), the 1-vs-2 boundary is marginally above chance (AUC 0.61--0.70), and all subsequent boundaries are at chance level (AUC $\leq$ 0.56).

\begin{table}[t]
\centering
\small
\begin{tabular}{lccccc}
\toprule
\textbf{Model} & \textbf{$K^*$} & \textbf{AUC} & \textbf{AUC} & \textbf{AUC} & \textbf{JSD} \\
 & & \textbf{0v1} & \textbf{1v2} & \textbf{2v3} & \textbf{0v1} \\
\midrule
Qwen2.5-1.5B & 2 & .996 & .70 & .59 & .267 \\
Pythia-1.4B   & 2 & .991 & .63 & .55 & .231 \\
OLMo-1B      & 2 & .988 & .61 & .54 & .218 \\
Qwen2.5-7B   & 2 & .994 & .65 & .56 & .258 \\
\bottomrule
\end{tabular}
\caption{Pairwise discrimination across model families and scales. $K^*=2$ at $\tau=0.60$ is universal. JSD decays by an order of magnitude from 0v1 to 1v2 in all cases.}
\label{tab:cross_model}
\end{table}

The JSD between consecutive generations exhibits rapid monotonic decay consistent with DPI predictions.
For Qwen2.5-1.5B, JSD drops from 0.267 (0v1) to 0.016 (1v2) to 0.007 (2v3)---a 38$\times$ reduction in two steps.

\paragraph{Scale does not help.}
Qwen2.5-7B ($5\times$ parameters) yields $K^*=2$ with 1v2 AUC of 0.65---\emph{lower} than the 1.5B variant (0.70).
Larger models produce higher-fidelity outputs, making consecutive generations \emph{more} similar and thus harder to distinguish, offsetting their greater capacity.

\subsection{Decoding Strategy Cannot Modulate $K^*$}

Table~\ref{tab:decoding} shows that $K^*=2$ is invariant across four decoding strategies spanning the deterministic-to-high-entropy spectrum.

\begin{table}[t]
\centering
\small
\begin{tabular}{lcccc}
\toprule
\textbf{Strategy} & \textbf{$K^*$} & \textbf{AUC 0v1} & \textbf{AUC 1v2} & \textbf{JSD 0v1} \\
\midrule
Greedy        & 2 & .997 & .62 & .305 \\
Nucleus $p$=0.9  & 2 & .996 & .68 & .280 \\
$T$=0.7       & 2 & .997 & .65 & .295 \\
$T$=1.2       & 2 & .993 & .64 & .240 \\
Nuc.\ $p$=0.95 (default) & 2 & .996 & .70 & .267 \\
\bottomrule
\end{tabular}
\caption{Decoding strategy ablation (Qwen2.5-1.5B). All strategies yield $K^*=2$. The 1v2 AUC range (0.62--0.70) is narrow and never approaches the utility threshold.}
\label{tab:decoding}
\end{table}

Greedy decoding maximizes the 0v1 JSD (0.305) because deterministic outputs are most distributional distinct from human text, but this does not propagate to deeper generations.
High temperature ($T$=1.2) reduces the 0v1 gap (JSD 0.240) while accelerating signal collapse at depth 2+.
No decoding configuration pushes $K^*$ to 3.

\subsection{$K^*$ Is Robust to Threshold Choice}

We verify that $K^*=2$ is not an artifact of the threshold $\tau=0.60$ by computing $K^*$ at four thresholds (Table~\ref{tab:threshold}).
At $\tau \geq 0.60$, all models yield $K^*=2$.
The sole exception---Qwen at $\tau=0.55$---reflects its 2v3 AUC of 0.59, marginally above threshold, within statistical noise.

\begin{table}[t]
\centering
\small
\begin{tabular}{lccc}
\toprule
\textbf{Threshold $\tau$} & \textbf{Qwen} & \textbf{Pythia} & \textbf{OLMo} \\
\midrule
0.55 & 3 & 2 & 2 \\
0.60 & 2 & 2 & 2 \\
0.65 & 2 & 2 & 2 \\
0.70 & 2 & 2 & 2 \\
\bottomrule
\end{tabular}
\caption{$K^*$ across AUC thresholds. Robustness reflects rapid signal collapse rather than measurement precision.}
\label{tab:threshold}
\end{table}

\subsection{Feature Representation Does Not Break the Ceiling}

Table~\ref{tab:features} compares handcrafted features with learned representations.
Frozen LM embeddings (4096-dim + MLP) achieve 1v2 AUC of 0.72 (+0.02 over handcrafted) and 2v3 AUC of 0.58, yielding the same $K^*=2$.
The ceiling is data-intrinsic, not a limitation of feature engineering.

\begin{table}[t]
\centering
\small
\begin{tabular}{lcccc}
\toprule
\textbf{Features} & \textbf{Acc.} & \textbf{MAE} & \textbf{AUC 1v2} & \textbf{$K^*$} \\
\midrule
Handcrafted (15D) & .37 & 1.08 & .70 & 2 \\
Learned (4096D) & .39 & 1.05 & .72 & 2 \\
\bottomrule
\end{tabular}
\caption{Handcrafted vs.\ learned representations. The 2pp accuracy gain from 4096-dim learned features does not shift $K^*$.}
\label{tab:features}
\end{table}

\subsection{Depth Estimation Accuracy Is Near-Chance}

Table~\ref{tab:classifiers} compares depth estimation methods.
CORAL achieves 37\% accuracy (6-class) with MAE 1.08---only 5pp above a binary classifier with random depth assignment.
All methods achieve near-random performance on depth 2--5 classes.

\begin{table}[t]
\centering
\small
\begin{tabular}{lccc}
\toprule
\textbf{Method} & \textbf{Accuracy} & \textbf{MAE} & \textbf{Macro-$F_1$} \\
\midrule
CORAL (ordinal) & .37 & 1.08 & .22 \\
Standard multiclass & .35 & 1.22 & .19 \\
Ordinal regression & .36 & 1.10 & .21 \\
Binary + random depth & .32 & 1.25 & .18 \\
\bottomrule
\end{tabular}
\caption{Classifier comparison. Ordinal structure provides minimal benefit when only two depth levels are meaningfully distinguishable.}
\label{tab:classifiers}
\end{table}

Frontier LLMs fare even worse: GPT-4o achieves 24\% depth accuracy (MAE 1.65) and Claude~3.5 achieves 22\% (MAE 1.72), both near the 20\% random baseline.
Notably, both LLMs retain reasonable binary accuracy (78--81\%), confirming that depth estimation is a fundamentally harder task than binary detection.

\subsection{Binary Detector Confidence Is Not a Depth Proxy}

Existing binary detectors' confidence scores carry negligible depth information.
DivEye and Binoculars achieve strong 0v1 AUC (0.91 and 0.94) but their 1v2 AUC drops to 0.54 and 0.56.
Spearman correlation between confidence and true depth is weak ($\rho = 0.18$ and $0.22$), and at depth 2+ the rank correlation is essentially zero ($\rho < 0.05$).

\subsection{Graduated Filtering Offers No Significant Advantage}

Table~\ref{tab:filtering} presents the downstream filtering evaluation.
Binary filtering improves over no filtering (+0.7pp MMLU, +0.9pp HellaSwag), confirming that removing synthetic text is beneficial.
However, graduated filtering provides only +0.1pp (MMLU) and +0.2pp (HellaSwag) over binary, both statistically non-significant ($p=0.62$ and $p=0.48$, respectively).

\begin{table}[t]
\centering
\small
\begin{tabular}{lcc}
\toprule
\textbf{Strategy} & \textbf{MMLU} & \textbf{HellaSwag} \\
\midrule
No filter & 25.1 \scriptsize{$\pm$0.2} & 42.3 \scriptsize{$\pm$0.3} \\
Binary & 25.8 \scriptsize{$\pm$0.3} & 43.2 \scriptsize{$\pm$0.4} \\
Graduated & 25.9 \scriptsize{$\pm$0.3} & 43.4 \scriptsize{$\pm$0.3} \\
Surprisal-based & 25.7 \scriptsize{$\pm$0.2} & 43.1 \scriptsize{$\pm$0.4} \\
\bottomrule
\end{tabular}
\caption{Downstream evaluation (3 seeds). Graduated filtering's advantage over binary is $\leq$0.2pp and non-significant ($p > 0.4$).}
\label{tab:filtering}
\end{table}

The dose-response experiment (Table~\ref{tab:dose}) confirms this null result scales with contamination severity: even at 50\% synthetic injection, graduated filtering gains at most 0.2pp over binary.

\begin{table}[t]
\centering
\small
\begin{tabular}{lcccc}
\toprule
\textbf{Dose} & \textbf{Bin.\ MMLU} & \textbf{Grad.\ MMLU} & \textbf{Bin.\ HS} & \textbf{Grad.\ HS} \\
\midrule
10\% & 25.6 & 25.6 & 43.0 & 43.0 \\
30\% & 25.5 & 25.6 & 42.8 & 42.9 \\
50\% & 25.2 & 25.3 & 42.3 & 42.5 \\
\bottomrule
\end{tabular}
\caption{Dose-response: graduated advantage does not scale with contamination level ($\Delta \leq 0.2$pp at all doses).}
\label{tab:dose}
\end{table}

\subsection{Cross-Model Transfer}

Table~\ref{tab:transfer} reveals a clean separation: binary detection transfers well across models (average accuracy 81\%), while depth-specific features do not (average 6-class accuracy 28\%, retention 40\% of in-domain performance).
Depth signal beyond the binary distinction is model-specific, requiring per-model classifier retraining---further limiting the practical utility of depth estimation.

\begin{table}[t]
\centering
\small
\begin{tabular}{lccc}
\toprule
\textbf{Train $\backslash$ Test} & \textbf{Qwen} & \textbf{Pythia} & \textbf{OLMo} \\
\midrule
\multicolumn{4}{c}{\emph{6-class accuracy}} \\
Qwen  & .37$^*$ & .28 & .26 \\
Pythia & .30 & .35$^*$ & .25 \\
OLMo  & .29 & .27 & .33$^*$ \\
\midrule
\multicolumn{4}{c}{\emph{Binary accuracy (human vs.\ synthetic)}} \\
Qwen  & .93$^*$ & .82 & .79 \\
Pythia & .84 & .91$^*$ & .77 \\
OLMo  & .83 & .80 & .88$^*$ \\
\bottomrule
\end{tabular}
\caption{Cross-model transfer. Binary detection transfers well (avg.\ 81\%); depth estimation does not (avg.\ 28\%). $^*$In-domain.}
\label{tab:transfer}
\end{table}

\section{Discussion}

\subsection{Why $K^*=2$? An Information-Theoretic Interpretation}

The DPI predicts monotonic decay of inter-generational mutual information, but does not specify the \emph{rate}.
Our measurements suggest that for 1--7B-scale models, the channel $T_\mathcal{M}$ is sufficiently lossy that a single pass eliminates the vast majority of distributional features distinguishing $P_0$ from $P_1$: the JSD drops by 17--23$\times$ from 0v1 to 1v2.
The second pass reduces the remaining signal to noise floor.

The invariance of $K^*$ across model families, scales (1B--7B), and decoding strategies implies that $K^*=2$ is not a property of any individual model but a regime-level characteristic of current-scale language models.
Whether larger models (70B+) would exhibit higher $K^*$ remains an open question beyond our computational scope.

\subsection{Implications for Data Governance}

Our results simplify the practical problem of synthetic text curation: \textbf{the actionable regime is binary}.
Practitioners should invest in robust first-generation detectors (human vs.\ synthetic) rather than depth estimators.
The three-class signal (human / gen-1 / gen-2+) is the finest granularity supported by current detection technology, and even this distinction provides no measurable downstream benefit over binary filtering.

\subsection{Limitations}

\paragraph{Prompt-chain vs.\ retrain-chain.}
Our primary experiments use prompt chains (model generates from prior outputs without retraining).
A retrain-chain validation with Pythia-1.4B (LoRA, rank 16) confirms $K^*=2$ but reveals statistically different feature dynamics (KS test $p=0.03$, JSD correlation $r=0.82$).
Prompt chains are thus a controlled proxy, not a faithful simulation of real-world recursive contamination.

\paragraph{Scale.}
Our 7B probe is limited to 1K documents and depth 0--3 due to GPU budget.
We cannot rule out that $K^*>2$ at significantly larger scales (e.g., 70B+).

\paragraph{Feature degradation ordering.}
Individual feature groups (surprisal, diversity, tail-thinning) show model-specific degradation orderings (cross-model Kendall $\tau = 0.33$), with all groups saturating at depth~2 regardless.
We do not claim a universal degradation ordering; this remains a descriptive observation.

\paragraph{Naturalistic validation.}
Applying our classifier to Common Crawl yields predicted depth--quality correlations of only $\rho=-0.15$ to $-0.19$, reflecting domain shift from controlled to in-the-wild data.

\subsection{Base vs.\ Instruct Models}

Instruction tuning (RLHF) increases 0v1 detectability (JSD 0.267 vs.\ 0.210 for Qwen2.5-1.5B Instruct vs.\ Base) by introducing distinctive stylistic patterns, but does not alter the recursive degradation dynamics: both variants yield $K^*=2$.
Base models may thus pose a greater challenge for binary detection in practice, while the depth estimation ceiling remains unchanged.

\section{Conclusion}

We present the first systematic measurement of the generational discrimination boundary $K^*$ for recursive synthetic text.
Across three model families, four decoding strategies, two model scales, and two feature representations, $K^*=2$ emerges as a hard, universal ceiling for 1--7B-parameter language models.
This is a direct empirical confirmation of the DPI prediction: the iterated lossy channel of recursive generation destroys inter-generational signal within two passes.

The practical consequence is stark: depth-aware data curation reduces to binary filtering.
Graduated filtering, ordinal classification, and depth estimation---despite their theoretical appeal---provide no statistically significant improvement over simply removing all detected synthetic text.
We urge the community to focus on robust first-generation detection and to treat depth estimation as a fundamentally limited task at current model scales.

\section*{Acknowledgments}
Omitted for review.

\bibliographystyle{plainnat}
\begin{thebibliography}{10}

\bibitem[Shumailov et~al.(2024)]{shumailov2024ai}
Ilia Shumailov, Zakhar Shumaylov, Yiren Zhao, Nicolas Papernot, Ross Anderson, and Yarin Gal.
\newblock {AI} models collapse when trained on recursively generated data.
\newblock \emph{Nature}, 631:755--759, 2024.

\bibitem[Mitchell et~al.(2023)]{mitchell2023detectgpt}
Eric Mitchell, Yoonho Lee, Alexander Khazatsky, Christopher~D Manning, and Chelsea Finn.
\newblock {DetectGPT}: Zero-shot machine-generated text detection using random perturbations.
\newblock In \emph{Proceedings of ICML}, 2023.

\bibitem[Hans et~al.(2024)]{hans2024binoculars}
Abhimanyu Hans, Avi Schwarzschild, Valeriia Cherepanova, Hamid Kazemi, Aniruddha Saha, Micah Goldblum, Jonas Geiping, and Tom Goldstein.
\newblock Spotting {LLMs} with {B}inoculars: Zero-shot detection of machine-generated text.
\newblock In \emph{Proceedings of ICML}, 2024.

\bibitem[Cao et~al.(2020)]{cao2020rank}
Wenzhi Cao, Vahid Mirjalili, and Sebastian Raschka.
\newblock Rank consistent ordinal regression for neural networks with a proportional odds model.
\newblock \emph{Pattern Recognition}, 108:107616, 2020.

\bibitem[Feng et~al.(2024)]{feng2024towards}
Jifan Feng, et~al.
\newblock Towards verifiable text generation with evolving memory.
\newblock \emph{arXiv preprint arXiv:2312.09075}, 2024.

\bibitem[Yi et~al.(2024)]{yi2024on}
Junjie Yi, et~al.
\newblock On the detectability of {LLM}-generated text.
\newblock \emph{arXiv preprint}, 2024.

\end{thebibliography}

\appendix

\section{Retrain-Chain Ecological Validity}
\label{app:retrain}

To assess ecological validity, we construct a retrain chain using Pythia-1.4B with LoRA fine-tuning (rank=16, $\alpha$=32) on C4 5K documents through depth 0--4.
Table~\ref{tab:retrain} compares prompt-chain and retrain-chain results.

\begin{table}[h]
\centering
\small
\begin{tabular}{lcc}
\toprule
\textbf{Metric} & \textbf{Prompt-chain} & \textbf{Retrain-chain} \\
\midrule
$K^*$ & 2 & 2 \\
AUC 0v1 & .991 & .972 \\
AUC 1v2 & .63 & .58 \\
AUC 2v3 & .55 & .53 \\
JSD 0v1 & .231 & .185 \\
JSD 1v2 & .012 & .008 \\
Feature KS-test $p$ & --- & .03 \\
\bottomrule
\end{tabular}
\caption{Prompt-chain vs.\ retrain-chain (Pythia-1.4B). $K^*$ is preserved; absolute signal levels are lower under retraining; feature distributions differ significantly.}
\label{tab:retrain}
\end{table}

$K^*=2$ is preserved under retraining, but absolute AUC and JSD values are uniformly lower, consistent with retraining accelerating the lossy-channel effect.
The significant KS test ($p=0.03$) indicates that while the saturation boundary is consistent, the feature-level dynamics differ, and prompt-chain classifiers should not be directly deployed on retrain-chain data without adaptation.

\section{Feature Group Degradation Analysis}
\label{app:features}

Table~\ref{tab:feat_group} presents per-feature-group 1v2 AUC by model.
All groups saturate at depth~2 regardless of model, but the \emph{ordering} of degradation (which group loses signal first) is model-specific (cross-model Kendall $\tau = 0.33$).

\begin{table}[h]
\centering
\small
\begin{tabular}{lccc}
\toprule
\textbf{Model} & \textbf{Surprisal} & \textbf{Diversity} & \textbf{Tail} \\
\midrule
Qwen2.5-1.5B & .63 & .58 & .61 \\
Pythia-1.4B & .56 & .60 & .59 \\
OLMo-1B & .58 & .59 & .55 \\
\bottomrule
\end{tabular}
\caption{Per-feature-group 1v2 AUC. All groups saturate at depth 2; degradation ordering is inconsistent across models.}
\label{tab:feat_group}
\end{table}

\section{LLM Zero-Shot Depth Estimation}
\label{app:llm}

\begin{table}[h]
\centering
\small
\begin{tabular}{lcccc}
\toprule
\textbf{Model} & \textbf{5-class Acc.} & \textbf{MAE} & \textbf{Binary Acc.} & \textbf{Kendall $\tau$} \\
\midrule
GPT-4o & .24 & 1.65 & .78 & .12 \\
Claude 3.5 & .22 & 1.72 & .81 & .10 \\
Random & .20 & 2.00 & .50 & .00 \\
CORAL (ours) & .37 & 1.08 & .93 & .35 \\
\bottomrule
\end{tabular}
\caption{LLM zero-shot depth estimation. Frontier LLMs achieve near-random depth accuracy despite reasonable binary detection.}
\label{tab:llm_zeroshot}
\end{table}

\section{Naturalistic Validation}
\label{app:naturalistic}

Applying our Qwen2.5-1.5B-trained classifier to Common Crawl 2025--2026 snapshots yields:
predicted 62\% human, 25\% gen-1, 13\% depth-2+.
Correlation of predicted depth with AlignScore is $-0.15$; with vocabulary richness, $-0.19$.
Binary synthetic prediction correlates more strongly with quality ($\rho=-0.28$), consistent with our main finding that depth signal beyond binary is weak.

\end{document}
